本章要回答的不是“资料够不够多”，而是哪些资料足以改变盈利、估值与动作。结论是：交易所披露足以支持合并口径财务、股本、聚合订单和情景估值；但逐船价格、排程、统一产能利用率、付款节点、军品经济性与可用Street目标价仍然缺失。缺失项不以行业常识补齐，而是通过合并敏感性、零计价或观察池处理。

\section{证据层级：事实优先，推导可复算，弱来源零权重}

同一个数字在不同来源中具有不同用途。公司年报可以证明上市公司聚合订单，却不能证明具体船厂的逐船毛利；行业数据库可以说明全球周期，却不能证明600150的订单归属；券商原始PDF可以说明市场预期，却不能替代公司实际披露。

\begin{exhibitbox}[证据金字塔与允许用途]
\small
\begin{tabularx}{\exhibitboxwidth}{L{2.3cm}L{3.7cm}L{4.0cm}X}
\toprule
\textbf{层级} & \textbf{主要来源} & \textbf{允许支持} & \textbf{不得越界} \\
\midrule
A1 & 上交所公司公告、审计年报、季报、重组文件、公司IR原文 & 法定口径、股本、聚合订单、交付、财务与主体边界 & 未披露合同字段、军品经济性、逐船毛利 \\
A2 & 工信部、IMO、OECD、中船协 & 国家/行业周期、政策状态、技术背景 & 600150公司份额、订单或EPS \\
B1 & 原始券商PDF、公开行业数据库 & 券商自身预测、评级及行业参考 & 公司实际、有效Street目标价（若原文未披露） \\
C/G & 第三方预览、媒体转述、失败探针、未取得原件 & 发现线索、记录证据缺口 & 估值锚、目标价、订单归属或可投资结论 \\
分析推导 & 以A1/B1输入可复算的模型与算术 & 情景收入、利润、EPS、隐含预期和公允价值 & 伪装成公司披露或逐合同事实 \\
\bottomrule
\end{tabularx}
\sourcenote{来源注册表；公司与市场主来源CF-03--CF-17、PR-01--03，行业来源II-06--II-12，券商来源BR-01--BR-13。}
\end{exhibitbox}

投资含义是：证据质量直接决定估值权重。13份原始券商PDF可用于盈利预测修订史，但都没有可用目标价，故Broker/Street目标价权重为0；第三方58.96元线索和未取得正文的47.00元线索同样为0。弱来源并未被删除，而是以“未披露/未取得”形式保留，避免读者误以为我们已穷尽并证实这些高目标价。

\section{事实、推导与禁止项}

读者需要区分三种句子。第一类是正式披露，例如2025年末民船/海工在手4,674.51亿元。第二类是透明推导，例如在手/2025收入约3.08倍。第三类是被阻断的推断，例如把在手订单按三年平均摊入收入。前两类可以进入报告，第三类必须留在证据缺口或敏感性中。

\begin{exhibitbox}[关键主张治理]
\small
\begin{tabularx}{\exhibitboxwidth}{L{4.0cm}L{2.4cm}L{3.6cm}X}
\toprule
\textbf{主张} & \textbf{证据状态} & \textbf{报告采用表述} & \textbf{估值处理} \\
\midrule
2025收入1,519.78亿元、归母78.48亿元 & A1审计披露 & 2025年同一控制下重述口径 & 作为合并基期 \\
2026H1归母92--110亿元 & A1正式预告，结果未审计 & 未经审计业绩预告，不机械年化 & 只作验证锚 \\
民船/海工在手4,674.51亿元 & A1聚合披露 & 合并后上市公司口径，不代表集团或完整军品 & 可见性，不直接计收入 \\
中远约500亿元/87艘项目 & A1归属、部分字段缺失 & 上市子公司承建，含意向且已在聚合订单内 & 去重质量证据，增量信用0 \\
沪东中华LNG能力 & 集团姐妹单位，上市联系缺失 & 集团能力背景，不属于七家主要船厂 & 订单、利润、注入价值均为0 \\
军品订单与利润 & 经济性未披露 & 能力背景，不做军品SOTP & 0 \\
弱来源目标价 & 原始材料不足 & 只记录缺口，不称Street目标价 & 0 \\
\bottomrule
\end{tabularx}
\sourcenote{CF-03、CF-05、CF-12；AG-03--AG-05；claim audit。}
\end{exhibitbox}

这套治理的“所以呢”是防止乐观因素被重复资本化：订单先进入交付和毛利模型，不再作为独立资产加回；集团能力不能下放给上市公司；未验证期权不进入熊、基、牛任一情景。这样做会牺牲部分故事弹性，却提高目标价可复算性。

\section{数据截止与剩余不确定性}

市场数据截至2026年7月22日盘后；融资融券最新日为7月21日，沪股通/基金为6月30日季度末存量。财务实际截至2026Q1，H1仅有未经审计预告。2025A与2024比较期均采用2025年同一控制下重述口径。当前估值使用75.2562亿股，2025A EPS保留公司披露的加权平均63.2964亿股分母。\srcid{CF-03, CF-04, CF-05, PR-01, PR-02, PR-03}

主要剩余缺口包括：逐船合同价与排程、船型级毛利、七家船厂统一利用率、付款里程碑、客户融资和撤单条款、军品订单经济性、有效券商目标价、严格自由流通分类以及日频沪股通逐证券净流。727日前复权K线、换手、两融、季度沪股通、基金和龙虎榜查询已闭环；剩余缺口不阻断合并估值，却阻断船型/船厂SOTP、军事业务溢价、严格自由流通推断、日频北向流判断和有效Street目标价加权。

\begin{sourcequalitybox}[证据边界对动作的约束]
报告可以给出合并口径公允价值，但不能声称掌握逐船盈利。公允价值与现价重合时，缺口应转化为“等待验证”动作，而不是靠更激进的故事填满上行空间。后续若取得正式中报、项目排程或原始目标价，只能通过重新审计模型进入结论，不能局部叠加。
\end{sourcequalitybox}

因此，本章的风险含义不是“数据不完美所以什么都不能判断”，而是精确指出可判断的层级：合并方向可信，细分持续期仍需验证；估值可做，期权不计；行动可以明确，但不应伪装成高置信度买入。
